\section{Phase 3: three real tasks, three precisions, three seeds}\label{sec:tiny}

Phases 1 and 2 used small grey pictures (MNIST and Fashion-MNIST), one training run each, and int8 only. That
is a good first test, but a reviewer would ask three fair questions:

\begin{enumerate}
  \item Does it still work on \emph{harder} tasks, and on things that are not pictures at all?
  \item Was the result luck? (One training run can be lucky or unlucky.)
  \item What happens with \emph{fewer} bits (int4) and \emph{more} bits (int16)?
\end{enumerate}

Phase 3 answers all three. It uses the tasks and the networks of \textbf{MLPerf Tiny}, the standard benchmark
for machine learning on tiny chips, trains each network \textbf{three times with different random seeds}, and
runs every network at \textbf{int4, int8 and int16} on the GPU \emph{and} on Ibex.

\begin{kid}
Imagine you invented a new way to do sums quickly in your head. Testing it on one page of easy homework is a
start. To be sure, you would test it on maths, on science and on music homework, try it on three different
days, and try it with small numbers and with huge numbers. That is exactly what this phase does.
\end{kid}

\subsection{The three tasks}

MLPerf Tiny picks one task from each kind of sensor a tiny device might have: a camera, a microphone and a
vibration or sound sensor on a machine. We use its reference network for each.

\begin{table}[H]
\centering
\small
\begin{tabularx}{\linewidth}{@{}l>{\raggedright\arraybackslash}X>{\raggedright\arraybackslash}X>{\raggedright\arraybackslash}X@{}}
\toprule
 & \textbf{Image classification} & \textbf{Keyword spotting} & \textbf{Anomaly detection}\\
\midrule
Dataset & CIFAR-10 & Google Speech Commands v2 & ToyADMOS (ToyCar, DCASE 2020)\\
What goes in & a $32\times32$ colour picture & one second of speech & a short slice of the sound of a toy car\\
What comes out & one of 10 objects (plane, cat, truck, \dots) & one of 12 answers: 10 words (``yes'', ``no'', ``up'', ``down'', \dots), ``silence'' or ``unknown word'' & ``normal'' or ``broken'' (an anomaly score)\\
Network & ResNet-8 & DS-CNN & fully connected autoencoder\\
Weights & 78{,}000 & 23{,}000 & 268{,}000\\
Multiplications per input & 12.5 million & 2.7 million & 264{,}000\\
Training / validation / test & 45{,}000 / 5{,}000 / 10{,}000 pictures & 38{,}461 / 4{,}627 / 4{,}890 clips & normal sounds only / 10\% held out / 2{,}459 clips (1{,}059 broken)\\
Score we report & top-1 accuracy & top-1 accuracy & AUC (explained below)\\
\bottomrule
\end{tabularx}
\caption{The three MLPerf Tiny tasks used here (\file{tiny/data.py}, \file{tiny/models.py}).}
\end{table}

\subsubsection*{ResNet-8 for pictures}

ResNet-8 is a small convolutional network with three ``stacks'' of 16, 32 and 64 filters. Its special trick is
the \emph{shortcut}: each stack adds its input back to its output (Figure~\ref{fig:resnet8}). Shortcuts make
deep networks much easier to train. When a stack makes the picture smaller (stride 2) or has more channels, the
shortcut goes through a small $1\times1$ convolution so the shapes match.

\begin{figure}[H]
\centering
\begin{tikzpicture}[node distance=0.35cm,
  blk/.style={draw,rounded corners,minimum height=0.75cm,minimum width=1.1cm,font=\scriptsize,align=center},
  >={Stealth[length=1.6mm]}]
  \node[blk,fill=gray!10] (in) {picture\\$32{\times}32{\times}3$};
  \node[blk,fill=blue!8,right=of in] (c1) {conv\\16};
  \node[blk,fill=blue!8,right=of c1] (a1) {conv\\16};
  \node[blk,fill=blue!8,right=of a1] (b1) {conv\\16};
  \node[circle,draw,inner sep=1pt,right=0.25cm of b1,font=\scriptsize] (p1) {$+$};
  \node[blk,fill=blue!14,right=0.3cm of p1] (a2) {conv\\32, /2};
  \node[blk,fill=blue!14,right=of a2] (b2) {conv\\32};
  \node[circle,draw,inner sep=1pt,right=0.25cm of b2,font=\scriptsize] (p2) {$+$};
  \node[blk,fill=blue!22,right=0.3cm of p2] (a3) {conv\\64, /2};
  \node[blk,fill=blue!22,right=of a3] (b3) {conv\\64};
  \node[circle,draw,inner sep=1pt,right=0.25cm of b3,font=\scriptsize] (p3) {$+$};
  \node[blk,fill=gray!10,below=0.6cm of p3] (out) {average +\\fc $\to$ 10};
  \draw[->] (in)--(c1); \draw[->] (c1)--(a1); \draw[->] (a1)--(b1); \draw[->] (b1)--(p1);
  \draw[->] (p1)--(a2); \draw[->] (a2)--(b2); \draw[->] (b2)--(p2);
  \draw[->] (p2)--(a3); \draw[->] (a3)--(b3); \draw[->] (b3)--(p3); \draw[->] (p3)--(out);
  \draw[->,orange!80!black] ($(c1.east)+(0.12,0)$) -- ++(0,0.75) -| node[pos=0.25,above,font=\tiny]{shortcut (same shape)} (p1);
  \draw[->,orange!80!black] ($(p1.east)+(0.12,0)$) -- ++(0,-0.75) -| node[pos=0.25,below,font=\tiny]{$1{\times}1$ conv, /2} (p2);
  \draw[->,orange!80!black] ($(p2.east)+(0.12,0)$) -- ++(0,0.75) -| node[pos=0.25,above,font=\tiny]{$1{\times}1$ conv, /2} (p3);
\end{tikzpicture}
\caption{ResNet-8. Every blue box is a $3\times3$ convolution with batch norm and ReLU. Orange: the shortcuts.}
\label{fig:resnet8}
\end{figure}

\subsubsection*{DS-CNN for spoken words}

A computer does not look at sound as a wiggly line. First it cuts the second of audio into 49 short pieces
(30 ms each, one every 20 ms) and, for each piece, measures how much energy there is at different pitches. That
gives a small ``picture of the sound'', $49\times10$ numbers, called \textbf{MFCC} features (mel-frequency
cepstral coefficients). From then on, keyword spotting is just classifying a small picture.

DS-CNN means \emph{depthwise-separable} CNN. A normal convolution mixes every input channel with every output
channel at every position, which is a lot of multiplications. DS-CNN splits that into two cheap steps:

\begin{itemize}
  \item a \textbf{depthwise} $3\times3$ convolution, which looks at neighbours but keeps each channel separate
        (9 multiplications per value), then
  \item a \textbf{pointwise} $1\times1$ convolution, which mixes channels but looks at only one position
        (64 multiplications per value).
\end{itemize}

\begin{example}[: Why depthwise-separable is cheaper]
For 64 channels in and 64 out, a normal $3\times3$ convolution needs $3\times3\times64 = 576$ multiplications
for every output value. The depthwise + pointwise pair needs $9 + 64 = 73$. That is about 8 times fewer, which
is why phones and watches use this design.
\end{example}

\subsubsection*{An autoencoder for broken machines}

Here there is no ``right answer'' to learn, because nobody records thousands of broken toy cars. Instead the
network learns what \emph{normal} sounds like. An \textbf{autoencoder} squeezes its input (640 numbers: 5 time
slices $\times$ 128 pitch bands) through a narrow middle of only 8 numbers and then tries to rebuild the
original 640. Trained on normal sounds only, it rebuilds normal sounds well. A broken car sounds different, the
network has never learned to rebuild that, so the rebuilt version is worse. The \emph{reconstruction error}
(how far the rebuilt sound is from the real one) is the anomaly score: big error means ``probably broken''.

\begin{kid}
Imagine you have heard your own front door close a thousand times. You could hum the sound perfectly. One day
it squeaks. When you try to hum what you just heard, your hum does not match, and that mismatch tells you
something is wrong. The autoencoder hums.
\end{kid}

\subsubsection*{What is AUC?}

For the broken-machine task we cannot say ``87\% correct'', because the answer is a score, not a class. We would
have to pick a threshold (``error above 12 means broken''), and every threshold gives a different accuracy.
\textbf{AUC} avoids picking one. It answers a simple question:

\begin{keyidea}
\textbf{AUC} = if you take one broken clip and one normal clip at random, how often does the broken one get the
higher anomaly score? 100\% is perfect, 50\% is no better than tossing a coin.
\end{keyidea}

\begin{example}[: AUC with 2 + 2 clips]
Normal clips get scores 3 and 5, broken clips get 4 and 9. There are $2\times2 = 4$ (broken, normal) pairs:
(4, 3) right, (4, 5) wrong, (9, 3) right, (9, 5) right. Three out of four are right, so AUC = 75\%.
\end{example}

We follow the DCASE 2020 rules: the score of a 10-second clip is the mean error over its slices, the AUC is
computed separately for each of the four toy cars in the test set, and we report the average of the four. The
official DCASE baseline for ToyCar reaches about 78.8\%, so our float network (77.7\%) is a faithful
reproduction.

\subsection{Three precisions: int4, int8 and int16}

Phase 2 only used int8. Phase 3 adds a smaller and a larger number size. Weights are signed and activations
are never negative (they come out of a ReLU), so they get slightly different ranges:

\begin{table}[H]
\centering
\small
\begin{tabular}{lrrl}
\toprule
 & weights & activations & how many different values\\
\midrule
int4 & $-7 \dots 7$ & $0 \dots 15$ & 15 and 16\\
int8 & $-127 \dots 127$ & $0 \dots 255$ & 255 and 256\\
int16 & $-32767 \dots 32767$ & $0 \dots 65535$ & 65{,}535 and 65{,}536\\
\bottomrule
\end{tabular}
\caption{The integer ranges (\code{QW} and \code{QA} in \file{tiny/qnet.py}).}
\end{table}

\begin{kid}
Think of a ruler. An int4 ruler has only 15 marks, so you have to round every length a lot. An int8 ruler has
255 marks. An int16 ruler has 65{,}535 marks, so fine that you can hardly see the rounding. Fewer marks means
smaller numbers (cheaper to store and multiply), but more rounding mistakes.
\end{kid}

The networks' inputs need a word too. CIFAR-10 pictures go in as raw pixels $0 \dots 255$ (at int4 they are
scaled to $0 \dots 15$), so black is 0, as on a real camera. Sound features (MFCC and the log-mel slices) can be
negative, so they are scaled into the signed weight range ($\pm7$, $\pm127$ or $\pm32767$).

\subsection{The methods we compare}

\begin{itemize}
  \item \textbf{Exact}: the shift-and-add multiply, equally optimised (fused loop, skips zeros, loops over the
        bits of the smaller number). This is the baseline everything is measured against.
  \item \textbf{Square table}: exact products from the polarization identity
        $xw = \big((x+w)^2 - x^2 - w^2\big)/2$ and a table of squares (Section~\ref{sec:ibexres}). Same answers
        as exact, so same accuracy.
  \item \textbf{RoBA-Dn1} and \textbf{RoBA-Dn2}: the approximate multipliers with the LOD instruction.
  \item \textbf{Approximate square} (int8 only): an experiment that combines the two ideas, see below.
\end{itemize}

\subsubsection*{Why the square table is ``not feasible'' at int16}

The table must hold $v^2$ for every value $x+w$ can take. That range is roughly from $-2\times$(largest weight)
to (largest activation $+$ largest weight):

\begin{center}
\small
\begin{tabular}{lrrl}
\toprule
 & entries & bytes & fits in a tiny chip?\\
\midrule
int4 & 37 & 148 B & easily\\
int8 & 637 & 2.5 KB & yes\\
int16 & 163{,}837 & 1.3 MB (squares need 64 bits) & no: more than the whole 1 MB memory\\
\bottomrule
\end{tabular}
\end{center}

So in every int16 row the square table is marked \emph{not feasible}. This is the exact situation where RoBA is
supposed to help: its cost does not grow with the table, because it has no table.

\subsubsection*{The approximate-square experiment}

If squares are the only multiplication left, could we make \emph{them} cheap with LOD instead of a table? A
square of $v$ with $k = \mathrm{LOD}(v)$ can be estimated as $\tilde S(v) = (2v - 2^k)\cdot 2^k$ (the same idea
as Mitchell's multiplier). This is only a few per cent too small for one square, but watch what happens when
three of them are combined:

\begin{example}[: Why approximate squares go wrong]
Take $x = 100$, $w = 3$, so the true product is 300.

\begin{tabular}{lrrr}
 & $v$ & true $v^2$ & $\tilde S(v)$, $k=\mathrm{LOD}(v)$\\
 & 103 & 10{,}609 & $(206 - 64)\cdot 64 = 9{,}088$\\
 & 100 & 10{,}000 & $(200 - 64)\cdot 64 = 8{,}704$\\
 & 3 & 9 & $(6 - 2)\cdot 2 = 8$\\
\end{tabular}

Exact: $(10609 - 10000 - 9)/2 = 300$. Approximate: $(9088 - 8704 - 8)/2 = \mathbf{188}$, which is 37\% too small.
The two big squares are each about 1{,}500 too small, but not by the \emph{same} amount, and the product is the
small difference between them. A small error on a big number becomes a huge error on the small answer.
\end{example}

\subsection{How the numbers were produced}

\begin{figure}[H]
\centering
\begin{tikzpicture}[node distance=0.45cm,
  st/.style={draw,rounded corners,fill=blue!5,minimum height=0.9cm,text width=2.35cm,align=center,font=\scriptsize},
  >={Stealth[length=1.6mm]}]
  \node[st] (tr) {\textbf{1. Train in float}\\3 seeds on the GPU};
  \node[st,right=of tr] (ca) {\textbf{2. Calibrate}\\pick the scales, try 4 percentiles};
  \node[st,right=of ca] (qa) {\textbf{3. Fine-tune (QAT)}\\with the integer and approximate arithmetic};
  \node[st,right=of qa] (ev) {\textbf{4. Evaluate}\\bit-exact integer model, full test set};
  \node[st,fill=orange!8,below=0.7cm of ev] (ib) {\textbf{5. Ibex}\\same integer program as C, seed 0, 20 inputs};
  \draw[->] (tr)--(ca); \draw[->] (ca)--(qa); \draw[->] (qa)--(ev); \draw[->] (ev)--(ib);
  \node[font=\scriptsize,left=0.3cm of ib,text width=6.6cm,align=right] {every output value on Ibex is checked to be
  \emph{identical} to the Python integer model};
\end{tikzpicture}
\caption{The Phase 3 pipeline (\file{tiny/run.py} for steps 1--4, \file{hw/run\_tiny.py} for step 5).}
\end{figure}

\paragraph{1. Training.} Each network is trained in normal floating point with the recipe of its reference:
ResNet-8 for 200 epochs and DS-CNN for 100 epochs with SGD, the autoencoder for 60 epochs with Adam. Pictures are
randomly shifted and flipped, sounds randomly shifted in time and mixed with background noise. This is done
three times with seeds 0, 1 and 2.

\begin{kid}
A \textbf{seed} is the starting number of the computer's dice. The network starts from random weights and sees
the training examples in a random order; a different seed means different dice rolls, so a slightly different
network at the end. If a result holds for all three seeds, it is not luck.
\end{kid}

\paragraph{2. Calibration.} To turn decimal numbers into integers we need a scale for each layer: which decimal
value becomes the largest integer (255 at int8)? Using the very largest value ever seen wastes most of the
ruler on rare outliers; cutting too low clips too many values. So we try four cut-offs (the 99.9th, 99.99th and
99.999th percentile, and the maximum) and keep the one with the best \emph{validation} accuracy. Batch norm is
folded into the weights first, as in Phase 2, and the weights get one scale per output channel.

\paragraph{3. Quantization-aware training (QAT).} Rounding to int4, or using a multiplier that is 9\% too small
on average, changes every layer's output a little. QAT lets the network get used to that: training continues
for a few epochs while the forward pass uses the integer rounding and the approximate multiplier, and the
backward pass pretends everything was exact (the straight-through estimator from Phase 1). We try two learning
rates and keep the best of ``no fine-tuning'', ``rate 1'' and ``rate 2'', again judged on the validation set
only. The approximate-square method is \emph{not} fine-tuned; it reuses the exact network.

\paragraph{4. Evaluation.} The test score comes from a Python model that does exactly the integer arithmetic the
chip does (64-bit sums, the same rounding, the same rescaling), over the \emph{whole} test set.

\paragraph{5. Ibex.} The same integer program is turned into C (\file{hw/run\_tiny.py} writes the weights and the
layer calls into \file{net.h}), compiled for RV32IC + LOD and run on the Verilator model of Ibex for seed 0,
with 20 inputs per configuration (for ToyADMOS: one slice from each of 20 clips). For all
$3 \text{ tasks} \times 12 \text{ configurations} = 36$ builds, every one of the output values (10 scores,
12 scores or 640 rebuilt numbers per input) is \emph{identical} to the Python model. So the test scores of step~4
are exactly the scores the chip would get, and the Ibex cycle counts belong to that very program.

\begin{careful}
The mean $\pm$ standard deviation below is over the 3 seeds. Cycle counts barely change from seed to seed (the
work is the same; only which activations happen to be zero changes), so we measured them for seed 0.
\end{careful}

\subsection{Accuracy}

\begin{table}[H]
\centering
\small
\begin{tabular}{llccc}
\toprule
 & & \textbf{CIFAR-10} & \textbf{Speech Commands} & \textbf{ToyADMOS}\\
 & & accuracy (\%) & accuracy (\%) & AUC (\%)\\
\midrule
float32 & & $87.38 \pm 0.12$ & $95.41 \pm 0.15$ & $77.68 \pm 0.41$\\
\midrule
\multirow{4}{*}{int4} & Exact = Square table & $85.45 \pm 0.14$ & $60.69 \pm 0.36$ & $77.13 \pm 0.12$\\
 & RoBA-Dn1 & $84.43 \pm 0.24$ & $61.04 \pm 0.44$ & $76.94 \pm 0.18$\\
 & RoBA-Dn2 & $85.49 \pm 0.03$ & $61.26 \pm 0.32$ & $77.39 \pm 0.10$\\
\midrule
\multirow{4}{*}{int8} & Exact = Square table & $87.40 \pm 0.04$ & $95.36 \pm 0.05$ & $77.86 \pm 0.33$\\
 & RoBA-Dn1 & $85.70 \pm 0.25$ & $94.70 \pm 0.34$ & $77.10 \pm 0.47$\\
 & RoBA-Dn2 & $87.29 \pm 0.03$ & $95.24 \pm 0.29$ & $77.73 \pm 0.56$\\
 & Approximate square & $23.39 \pm 4.27$ & $27.88 \pm 11.31$ & $54.13 \pm 6.99$\\
\midrule
\multirow{4}{*}{int16} & Exact & $87.38 \pm 0.06$ & $95.47 \pm 0.12$ & $77.68 \pm 0.41$\\
 & Square table & \multicolumn{3}{c}{\emph{not feasible (table larger than the memory)}}\\
 & RoBA-Dn1 & $85.46 \pm 0.12$ & $94.56 \pm 0.09$ & $77.22 \pm 0.32$\\
 & RoBA-Dn2 & $87.26 \pm 0.11$ & $95.29 \pm 0.09$ & $78.04 \pm 0.17$\\
\bottomrule
\end{tabular}
\caption{Test score, mean $\pm$ standard deviation over 3 seeds, full test sets
(\file{results/tiny/<dataset>/<model>\_accuracy.csv}). The square table gives exactly the same products as
exact, so its accuracy is the same.}
\label{tab:tinyacc}
\end{table}

The differences that matter are against \emph{exact at the same precision}, because that separates the
effect of the multiplier from the effect of rounding to fewer bits. Figure~\ref{fig:tinyacc} shows them. Each
bar is computed seed by seed (Dn1 with seed 0 minus exact with seed 0, and so on), then averaged, which cancels
most of the seed-to-seed noise.

\begin{figure}[H]
\centering
\includegraphics[width=\linewidth]{tiny_acc}
\caption{Accuracy (AUC for ToyADMOS) of RoBA minus exact at the same precision, in percentage points.
Bars: mean over 3 seeds; whiskers: standard deviation.}
\label{fig:tinyacc}
\end{figure}

What the table and the figure say:

\begin{enumerate}
  \item \textbf{int8 and int16 are as good as float.} With the exact multiplier, int8 is within 0.2 points of
        float on all three tasks, and int16 is the same as float. So any loss below comes from the multiplier,
        not from the integers.
  \item \textbf{RoBA-Dn2 is practically free.} At int8 and int16 it is never more than 0.2 points below exact
        (and on ToyADMOS int16 it is 0.36 points \emph{above}, which is noise in the lucky direction).
  \item \textbf{RoBA-Dn1 costs a little.} About 1.7--1.9 points on CIFAR-10, the hardest task, and 0.5--0.9 points
        on the two sound tasks. Dn1 always rounds \emph{down}, so its products are 9\% too small on average, and
        fine-tuning can only partly hide that.
  \item \textbf{int4 costs about 2 points on CIFAR-10} already with the exact multiplier. On top of that, Dn2
        costs nothing and Dn1 about one more point. On the two sound tasks the int4 differences between methods
        are smaller than the seed-to-seed spread.
  \item \textbf{The approximate square fails completely} (23\%, 28\% and 54\%, close to guessing), exactly as
        the example above predicted. It is a useful negative result: the polarization identity needs
        \emph{exact} squares.
\end{enumerate}

\begin{careful}
\textbf{Keyword spotting at int4 does not work yet} (about 61\% instead of 95\%), and this is the same for every
method, including exact. So it is a problem of our int4 recipe for this network, not of RoBA. Before
fine-tuning the int4 network scored only about 27\% on the validation set. Our best guess is the input: the
MFCC features are squeezed into just 15 levels ($-7 \dots 7$), and the first layer has only 40 weights per
filter, also at 4 bits. Low-bit papers usually keep the first layer at 8 bits. We have not tested that yet, so
the int4 keyword-spotting numbers are only useful for the \emph{relative} comparison of methods.
\end{careful}

\subsection{Speed on Ibex}

\begin{table}[H]
\centering
\small
\begin{tabular}{llrrr}
\toprule
 & & \textbf{CIFAR-10} & \textbf{Speech Commands} & \textbf{ToyADMOS}\\
\midrule
\multirow{4}{*}{int4} & Exact & 259.5 & 59.96 & 6.50\\
 & Square table & 185.0 \;(1.40$\times$) & 52.99 \;(1.13$\times$) & 4.92 \;(1.32$\times$)\\
 & RoBA-Dn1 & 224.9 \;(1.15$\times$) & 49.52 \;(1.21$\times$) & 5.95 \;(1.09$\times$)\\
 & RoBA-Dn2 & 252.1 \;(1.03$\times$) & 55.14 \;(1.09$\times$) & 6.74 \;(0.96$\times$)\\
\midrule
\multirow{5}{*}{int8} & Exact & 469.2 & 123.49 & 13.28\\
 & Square table & \textbf{197.9 \;(2.37$\times$)} & \textbf{62.83 \;(1.97$\times$)} & \textbf{5.12 \;(2.59$\times$)}\\
 & RoBA-Dn1 & 264.9 \;(1.77$\times$) & 70.34 \;(1.76$\times$) & 6.91 \;(1.92$\times$)\\
 & RoBA-Dn2 & 331.7 \;(1.41$\times$) & 85.52 \;(1.44$\times$) & 9.11 \;(1.46$\times$)\\
 & Approximate square & 288.8 \;(1.62$\times$) & 82.20 \;(1.50$\times$) & 7.21 \;(1.84$\times$)\\
\midrule
\multirow{4}{*}{int16} & Exact & 938.4 & 257.18 & 26.71\\
 & Square table & \multicolumn{3}{c}{\emph{not feasible}}\\
 & RoBA-Dn1 & \textbf{308.7 \;(3.04$\times$)} & \textbf{87.55 \;(2.94$\times$)} & \textbf{8.02 \;(3.33$\times$)}\\
 & RoBA-Dn2 & 393.3 \;(2.39$\times$) & 106.59 \;(2.41$\times$) & 10.58 \;(2.52$\times$)\\
\bottomrule
\end{tabular}
\caption{Million clock cycles per input on Ibex (seed 0, mean of 20 inputs), and in brackets the speed-up over
exact at the same precision (\file{results/tiny/<dataset>/<model>\_ibex\_summary.csv}). Bold: the fastest
method at each precision.}
\label{tab:tinyibex}
\end{table}

\begin{figure}[H]
\centering
\includegraphics[width=\linewidth]{tiny_speed}
\caption{Speed-up over exact at the same precision. The dashed line is ``as fast as exact''.}
\label{fig:tinyspeed}
\end{figure}

The same picture appears on all three tasks, even though they are a picture, a sound and a machine:

\begin{keyidea}
\textbf{The winner depends on the number size.}
\begin{itemize}
  \item \textbf{int4}: everything is about as fast as exact (1.0--1.4$\times$). Small numbers are already cheap
        to multiply.
  \item \textbf{int8}: the \textbf{square table} wins (2.0--2.6$\times$), with no accuracy loss at all.
  \item \textbf{int16}: the table no longer fits and \textbf{RoBA-Dn1} wins (2.9--3.3$\times$); RoBA-Dn2 still gives
        2.4--2.5$\times$ with essentially exact accuracy.
\end{itemize}
\end{keyidea}

\subsubsection*{Why the winner changes: the cost of one product}

Figure~\ref{fig:tinycycles} shows the same cycle counts as lines against the number size. The exact line
roughly \emph{doubles} from int4 to int8 and again from int8 to int16. The RoBA lines barely rise.

\begin{figure}[H]
\centering
\includegraphics[width=\linewidth]{tiny_cycles}
\caption{Million cycles per input against precision. Exact grows with the number of bits; RoBA hardly does;
the square table stops at int8.}
\label{fig:tinycycles}
\end{figure}

\begin{itemize}
  \item \textbf{Exact} shift-and-add does one step for every bit of the smaller number. 4 bits: about 4 steps.
        16 bits: about 16 steps. Its cost grows with the precision.
  \item \textbf{RoBA} finds the leading one with a single LOD instruction and then does a fixed handful of shifts
        and adds (Dn1: three terms; Dn2: twice that). Its cost is \emph{almost the same at any precision}.
        The small rise comes from the parts that are not multiplications: loading bigger weights, 64-bit sums at
        int16, and rescaling.
  \item \textbf{The square table} costs one memory read per product, whatever the precision, which is the cheapest
        of all, until the table does not fit.
\end{itemize}

\begin{example}[: Crossing lines, CIFAR-10]
Exact needs 259.5 million cycles at int4 and 938.4 million at int16: 3.6 times more work for the same network.
RoBA-Dn1 goes from 224.9 to 308.7 million: only 1.4 times more. At int4 the two are 1.15$\times$ apart; at int16
they are 3.04$\times$ apart. The more bits, the more RoBA saves.
\end{example}

At int4, RoBA-Dn2 on ToyADMOS is even a little \emph{slower} than exact (0.96$\times$): when the numbers have only
3 bits of size, the exact loop runs at most 3 times, and Dn2's two rounds of work cost more than that.

\subsubsection*{Where the cycles go inside ResNet-8}

\begin{figure}[H]
\centering
\includegraphics[width=\linewidth]{tiny_layers}
\caption{Cycles per layer for one CIFAR-10 picture at int8. conv1 is the first layer, s1x/s2x/s3x the three
stacks, s2s and s3s the $1\times1$ shortcut convolutions.}
\label{fig:tinylayers}
\end{figure}

The square table's lead is huge in the first layers (conv1: 3.4$\times$, s1a: 3.2$\times$) and small in the last
ones (s3b: 1.3$\times$, the same as Dn1). The reason is \textbf{zeros}. After each ReLU, many activations are 0,
and deeper layers have more of them. The exact and RoBA loops check for zero and skip those products for free;
the table loop does one table read per product without checking. So the deeper the layer, the more the other
methods catch up. This also means the table's speed-up would be smaller still on a network with more zeros.

\subsection{Energy}

Energy is estimated exactly as in Phase 2: cycles times the core's relative power. The core with the LOD
instruction uses 4.5\% more power ($0.829$ instead of $0.793$, relative to a core with a hardware multiplier),
so RoBA and the approximate square pay that extra; exact and the square table do not use LOD and do not pay it.

\begin{example}[: Energy of Dn1 at int16, CIFAR-10]
$\dfrac{308.7 \times 0.829}{938.4 \times 0.793} = \dfrac{255.9}{744.1} = 0.34$. RoBA-Dn1 uses about a third of the
energy of exact for one picture.
\end{example}

\begin{table}[H]
\centering
\small
\begin{tabular}{llccc}
\toprule
 & & CIFAR-10 & Speech Commands & ToyADMOS\\
\midrule
\multirow{3}{*}{int4} & Square table & 0.71 & 0.88 & 0.76\\
 & RoBA-Dn1 & 0.91 & 0.86 & 0.96\\
 & RoBA-Dn2 & 1.02 & 0.96 & 1.08\\
\midrule
\multirow{4}{*}{int8} & Square table & \textbf{0.42} & \textbf{0.51} & \textbf{0.39}\\
 & RoBA-Dn1 & 0.59 & 0.60 & 0.54\\
 & RoBA-Dn2 & 0.74 & 0.72 & 0.72\\
 & Approximate square & 0.64 & 0.70 & 0.57\\
\midrule
\multirow{2}{*}{int16} & RoBA-Dn1 & \textbf{0.34} & \textbf{0.36} & \textbf{0.31}\\
 & RoBA-Dn2 & 0.44 & 0.43 & 0.41\\
\bottomrule
\end{tabular}
\caption{Energy per input relative to exact at the same precision (exact = 1). Below 1 saves energy.}
\end{table}

Energy follows the speed-ups almost one to one. The only cases above 1 are RoBA-Dn2 at int4 on CIFAR-10 and
ToyADMOS, where the 4.5\% extra power is not paid back by the tiny speed-up.

\subsection{The trade-off in one picture}

\begin{figure}[H]
\centering
\includegraphics[width=0.85\linewidth]{tiny_tradeoff}
\caption{Every task, precision and method in one plot: speed-up over exact (right is better) against accuracy
change (up is better). Colour: method. Shape: precision. The approximate square is left out (it would sit far
below the plot).}
\label{fig:tinytradeoff}
\end{figure}

The best place to be is the top right. Three groups stand out:
\begin{itemize}
  \item the purple circles (square table, int8) sit right on the zero line at 2--2.6$\times$: free speed;
  \item the orange triangles (RoBA-Dn2, int16) sit on the zero line at about 2.4--2.5$\times$: almost free speed
        where the table cannot go;
  \item the blue triangles (RoBA-Dn1, int16) go furthest right (about 3$\times$) but up to 2 points lower: the
        choice when speed and energy matter more than the last point of accuracy.
\end{itemize}

\subsection{What Phase 3 adds to the story}

\begin{enumerate}
  \item \textbf{The results hold across domains.} A picture task, a speech task and a machine-sound task give
        the same ranking of methods at every precision.
  \item \textbf{The results are not luck.} With three seeds, the spread is small (usually 0.1--0.4 points), much
        smaller than the effects we report.
  \item \textbf{Precision decides the winner.} At int8 use the square table; at int16 use RoBA (Dn2 when accuracy
        matters, Dn1 when energy matters); at int4 none of them is worth much.
  \item \textbf{RoBA's real advantage is that its cost does not grow with the number of bits}, thanks to the
        single-cycle LOD instruction. The bigger the numbers, the bigger the win: up to 3.3$\times$ faster and
        about a third of the energy at int16.
  \item \textbf{Approximate squares do not work}: the subtraction in the polarization identity magnifies their
        error.
  \item \textbf{Every Ibex result is bit-exact} with the Python integer model (36 builds, all outputs identical).
\end{enumerate}
